\section{证据层级}
本报告把公司公告、定期报告、审计报告、实时行情、龙虎榜和媒体转述严格分层。官方年报和Q1报告用于历史财务、产品结构和股东表；7月14日H1预告用于最新盈利分母；公开行情用于价格和流动性；龙虎榜用于事件资金结构。搜索片段、雪球、财富号以及未经公司确认的客户和订单表述不进入基础估值。

\begin{exhibitbox}[表2-1\quad 已确认事实与不能越过的边界]
\begin{tabularx}{\exhibitboxwidth}{L{3.0cm}X L{3.0cm}X}
\textbf{事实} & \textbf{确认内容} & \textbf{不能推出} & \textbf{模型处理} \\
H1预告 & 归母1.70--2.40亿元；扣非2.30--3.30亿元；EPS0.1382--0.1951元 & 未经审计、没有分部利润 & 作为2026E硬分母范围 \\
Q1业绩 & 收入10.64亿元；归母0.51亿元；扣非1.03亿元 & 不能直接年化且归母受公允价值影响 & 用H1预告校准H2桥 \\
光纤销量 & 2025年销售300.48万芯公里，同比+4061.77\% & 不能推出2026年ASP、利用率和利润率 & 周期修复证据，不给成长倍数 \\
业务结构 & 电力电缆38.27\%、通信电缆26.58\%、安全业务19.61\% & 不能把公司称为纯光纤/光模块 & 综合平台PE/PB/PS \\
资金结构 & Q1基金/HKSCC/高盛等持仓；龙虎榜机构多空交替 & 不能推断7月具体基金在加仓 & 机构参与+主动交易，不作实名推断 \\
\end{tabularx}
\sourcenote{data/claim\_audit.md；data/source\_registry.md}
\end{exhibitbox}

\section{公司结构与财务底盘}
通鼎互联在2025年实现营业收入CNY34.13亿元，但归母净利润亏损CNY0.76亿元、扣非净利润亏损CNY0.41亿元。业务由光纤光缆、通信电缆、电力电缆、通信设备、安全业务和新能源组成。2026Q1收入升至CNY10.64亿元，扣非净利润CNY1.03亿元，但经营活动现金流仍为净流出CNY0.18亿元，说明利润修复尚未完全转化为现金。

\begin{exhibitbox}[表2-2\quad 2025年产品收入与毛利率]
\begin{tabularx}{\exhibitboxwidth}{L{3.2cm}R{2.2cm}R{1.8cm}R{2.0cm}X}
\textbf{产品} & \textbf{收入CNY100m} & \textbf{占比} & \textbf{毛利率} & \textbf{含义} \\
光纤光缆 & 1.83 & 5.34\% & 18.10\% & 周期修复弹性，但收入纯度有限 \\
通信电缆 & 9.07 & 26.58\% & 24.32\% & 传统通信线缆利润池 \\
电力电缆 & 13.06 & 38.27\% & 12.26\% & 最大收入底盘，毛利较低 \\
通信设备 & 2.59 & 7.58\% & 13.00\% & 辅助设备业务 \\
安全业务 & 6.69 & 19.61\% & 29.25\% & 第二利润池，仍需看并表后现金质量 \\
\end{tabularx}
\sourcenote{2025年年度报告，分产品收入与毛利率表；数据单位为CNY100m。}
\end{exhibitbox}

\section{资产负债表与会计风险}
2026Q1应收账款CNY14.51亿元、存货CNY11.79亿元，短期借款CNY22.34亿元，现金CNY9.51亿元，归母权益CNY25.20亿元。公司年报把收入确认、应收账款和存货跌价列为关键审计事项；这不是对报表的否定，但意味着在利润快速修复阶段，回款、账龄和库存可变现价值必须进入估值跟踪。
